\section{Abordagem Proposta}
\label{sec:metodo}

A abordagem proposta tem três componentes, todos construídos sobre a segmentação da transcrição em
turnos de fala descrita na Seção~\ref{sec:dados-particao}. O primeiro, a Unidade Deliberativa
Verificável (UDV), associa cada uma das 2.203 opiniões do resumo estruturado a um trecho da fala da
própria pessoa na transcrição e registra a posição exata desse trecho, o tipo de vínculo encontrado e
a origem do vínculo. O segundo reúne as falas de cada pessoa em todas as audiências de que ela
participou e usa um LLM para gerar, somente a partir dessas falas, um perfil textual da pessoa. O
terceiro fornece esse perfil, com trechos das falas anteriores da pessoa, a um LLM que deve
responder como ela responderia, e usa as UDVs de audiências que não entraram no perfil como gabarito
dessa simulação.

\subsection{Visão geral da arquitetura}
\label{sec:metodo-visao}

\input{figures/architecture}

No fluxo da Figura~\ref{fig:arquitetura}, cada transcrição é dividida em turnos, o cabeçalho
de cada turno identifica o autor do trecho, e o processamento segue a partir daí em dois ramos, que
se reencontram na simulação. No ramo da UDV (Seção~\ref{sec:metodo-udv}), cada audiência é tratada
isoladamente: os participantes citados na matéria são associados aos seus turnos, e cada opinião
recebe como evidência uma sentença da fala da pessoa, encontrada por citação literal ou por
similaridade semântica, com offsets de caracteres e um nível que declara o tipo de vínculo. No ramo
dos atores (Seções~\ref{sec:metodo-atores} e~\ref{sec:metodo-perfis}), os turnos são agrupados por
pessoa entre audiências, filtrados, restritos às audiências de treino e fornecidos a um LLM, que
escreve o perfil. Na simulação (Seção~\ref{sec:metodo-simulacao}), o perfil e os trechos
recuperados das falas de treino da pessoa compõem o prompt de um LLM, e as UDVs das audiências de
validação e de teste definem as perguntas com que a simulação é avaliada (Seção~\ref{sec:setup}).

Os dois ramos resolvem a identidade das pessoas de formas diferentes: a UDV associa o nome citado na
matéria aos cabeçalhos de uma única audiência, e o ramo dos atores associa entre si cabeçalhos de
audiências diferentes. Como ambos usam as mesmas chaves (audiência, índice do turno e offsets na
transcrição), o turno que contém a evidência de uma UDV identifica o registro de ator
correspondente, e é por esse cruzamento que uma opinião publicada sobre uma audiência de teste se
torna uma pergunta sobre um ator com perfil.

\subsection{Unidades Deliberativas Verificáveis}
\label{sec:metodo-udv}

As opiniões do PublicHearingBR registram o que a matéria atribui a cada pessoa, mas o conjunto não
indica de que ponto da transcrição cada opinião foi extraída (Seção~\ref{sec:dados}). Sem essa
indicação, não é possível apresentar ao leitor a passagem que originou uma opinião, verificar se a
fala a sustenta nem analisar as posições de uma pessoa a partir do que ela disse. A UDV estabelece
esse vínculo com um registro por opinião que aponta o trecho da transcrição usado como evidência
(Figura~\ref{fig:udv}).

\textbf{Resolução da pessoa.} O resumo estruturado registra o nome usado na matéria, que costuma diferir do
cabeçalho da transcrição (``Rodrigo Marinho'' na matéria, ``RODRIGO SARAIVA MARINHO'' na
transcrição), e quem preside a sessão aparece como \textit{PRESIDENTE}, com o nome entre parênteses.
Cada turno fornece nomes candidatos: o nome do cabeçalho e, nos turnos de presidência, também o nome
entre parênteses. Um participante é associado a um candidato quando os dois conjuntos de palavras,
sem acento e em maiúsculas, são iguais ou um contém o outro, desde que ambos tenham ao menos duas
palavras. Um nome formado por uma palavra é associado a um candidato idêntico ou, na falta deste, ao
único falante da audiência que tenha essa palavra no nome. Com essa regra, 1.020 dos 1.065
participantes são associados a ao menos um turno, e os 45 restantes respondem por 90 opiniões.

\input{figures/udv-example}

\textbf{Sentenças candidatas.} A segmentação em sentenças é feita dentro de cada turno da pessoa, e
cada sentença guarda o índice do turno de origem. Com isso, nenhuma sentença junta o fim de um turno
ao início do turno seguinte da mesma pessoa, entre os quais a transcrição registra falas de outros
participantes. São mantidas as sentenças com quatro palavras ou mais que não se reduzam a
marcação de palco, como ``(Palmas.)'', o que resulta em 115.599 sentenças candidatas nas 206
audiências.

\textbf{Citação literal.} Muitas opiniões trazem trechos entre aspas, que a matéria apresenta como
fala literal e que permitem um vínculo sem uso de modelo. De cada citação com 10 caracteres ou mais
são extraídos prefixos de 10, 6, 4 e 3 palavras, buscados nessa ordem em cada turno da pessoa, sem
distinção entre maiúsculas e minúsculas e respeitando as fronteiras de palavra; entre as citações
de uma mesma opinião, prevalece a de prefixo encontrado mais longo. Quando esse prefixo tem seis palavras
ou mais, a sentença que o contém é a evidência, com o nível \texttt{quote\_found}. Se o prefixo
ocorre mais de uma vez na fala, é escolhida a sentença com maior índice de Jaccard de palavras em
relação à opinião. O nível não exige que o restante da citação coincida com a fala: na
Figura~\ref{fig:udv}, a matéria escreve ``presidente da comissão, que é eleita'', e a transcrição,
``presidente que foi eleita''.

Prefixos de três ou quatro palavras costumam ser expressões comuns, encontradas em pontos da fala sem
relação com a citação. O limite de seis palavras foi definido a partir de uma leitura exploratória
dos casamentos agrupados por tamanho de prefixo, sem protocolo de anotação. Das 2.203 opiniões, 960
têm alguma citação, 549 têm algum prefixo encontrado na fala da pessoa e 277 têm prefixo de seis
palavras ou mais.

\textbf{Similaridade semântica.} Quando a opinião não tem citação confiável, ela e as sentenças
candidatas da pessoa são codificadas pelo Serafim 335m~\citep{gomes2024serafim}, um codificador de
sentenças para o português (vetores de 1.024 dimensões, entrada limitada a 128 tokens), e a
evidência é a sentença de maior similaridade de cosseno com a opinião. Um codificador denso foi
preferido ao TF-IDF porque a matéria costuma parafrasear a fala, e a comparação por vocabulário
compartilhado falha nesses casos. Se um prefixo curto (menos de seis palavras) de alguma citação da
opinião ocorre na sentença escolhida, o tipo de suporte registra a coincidência como corroboração
fraca, sem alterar o nível. Cada UDV registra a versão do codificador e o corte com que foi
produzida.

\textbf{Nível e proveniência.} O nível resume o tipo de vínculo encontrado
(Tabela~\ref{tab:udv-niveis}), e os níveis \texttt{semantic\_match\_high} e
\texttt{semantic\_match\_weak} indicam apenas de que lado do corte ficou o cosseno da melhor
sentença, sem afirmar que ela sustenta a opinião. A proveniência registra a origem do vínculo:
\texttt{weak} (supervisão fraca, isto é, rótulo produzido por regra, sem conferência humana) para o
casamento de texto e \texttt{model} para o codificador. Nenhum registro tem proveniência humana.

\textbf{Localização e cobertura.} A sentença escolhida é localizada na transcrição original por uma
expressão regular restrita ao turno de origem, e todas as 2.105 evidências foram localizadas com
offsets de caracteres. As 98 opiniões restantes, nos níveis \texttt{no\_evidence} e
\texttt{person\_not\_resolved}, não têm evidência; esses níveis registram que a fala da pessoa não
foi encontrada e não indicam erro da matéria. Os 8 registros \texttt{no\_evidence} são de
participantes cuja única fala transcrita é a marcação ``(Manifestação em LIBRAS.)'', e entre as
pessoas não resolvidas há casos de grafia divergente do nome, instituições citadas como
participantes e pessoas que não falaram na sessão.

\input{tables/udv-tiers}

\textbf{Calibração do corte.} O corte que separa \texttt{semantic\_match\_high} de
\texttt{semantic\_match\_weak} foi calibrado apenas nas audiências de treino, porque a amostra de
validação humana é sorteada no teste e depende dessa fronteira. A regra definida antes de qualquer
resultado escolheria o corte que maximizasse o índice J de Youden~\citep{youden1950index} na separação entre acertos e erros
do codificador em opiniões das quais os trechos entre aspas foram removidos. Há acerto quando a
sentença mais similar à opinião é a da própria citação. Essa regra se mostrou
degenerada: só 12 das 142 consultas do treino são acertos, os erros têm cosseno mais alto que os
acertos (medianas de 0,655 e 0,578, respectivamente), e o maior valor de J é zero, obtido por um
corte que aceita todas as consultas.

O corte adotado usa como positivos as 183 opiniões do treino com citação confiável, cada uma pareada
com a sentença da própria citação, e como negativos a mesma opinião pareada com uma sentença sorteada
de outra audiência do treino. Seu valor é o ponto médio entre o primeiro quartil dos positivos
(0,628) e o terceiro quartil dos negativos (0,281), igual a 0,454 e arredondado para 0,45 (IC 95\%
por \textit{bootstrap} de audiências, de 0,437 a 0,469). Com negativos sorteados da fala da própria
pessoa, que corresponde à escolha que o codificador de fato realiza, a mesma regra resultaria em
0,50. O corte define uma fronteira entre rótulos e não foi validado como medida de confiança.

\textbf{Verificação.} Um verificador separado refaz, com as mesmas funções e a partir das transcrições e dos resumos
estruturados, a
resolução de pessoa, a extração das sentenças candidatas e a busca de citações de cada opinião. Em
seguida, confere se cada registro respeita as regras: opinião e ator iguais aos do resumo estruturado, texto da
evidência contido em um único turno do ator, texto recuperado pelos offsets igual ao registrado e
nível coerente com o tipo de suporte, o cosseno e o corte. O verificador não recalcula a escolha do
codificador nem avalia se a evidência sustenta a opinião. O conjunto de UDVs descrito aqui não
apresenta nenhuma inconsistência, e uma versão anterior do conjunto, regenerada com o código que a
produziu, foi idêntica byte a byte à original.

\subsection{Falas por ator entre audiências}
\label{sec:metodo-atores}

Como a UDV trata uma audiência de cada vez, descrever o que uma pessoa defende ao longo do tempo
exige reunir suas falas em todas as audiências de que participou e, para isso, decidir quando
cabeçalhos de audiências diferentes se referem à mesma pessoa. A regra de nome contido usada na UDV
é aceitável dentro de uma audiência, que tem poucos falantes, mas entre audiências ela une pessoas
diferentes, como ``Luiz Lima'' e ``Gustavo Luiz de Lima Correia''. Por isso, o registro de ator usa
como chave o nome do cabeçalho sem acento e em maiúsculas (nos turnos de presidência, o nome entre
parênteses), complementado por uma lista de mesclas.

A lista de mesclas resulta da revisão, um a um, dos 153 pares de nomes candidatos (pares em que um
nome contém o outro ou os dois têm as mesmas palavras). Cada par foi examinado com
evidências do próprio conjunto: o cargo informado na matéria, a apresentação feita pelo presidente
da sessão, a autoapresentação, o partido e a UF do cabeçalho e a data da audiência. A revisão
produziu 37 grupos de mescla, uma chave dividida por audiência (um caso de homônimo exato) e um par
sem resolução, mantido separado. Nenhuma identidade foi conferida em registros externos ao conjunto.

Três regras de filtragem retiram os turnos que não trazem a fala da própria pessoa ou que se limitam
à condução da sessão: saem os turnos de intérprete e de tradução simultânea, cuja fala pertence a
outro participante, os turnos de presidência com menos de 50 palavras e os turnos vazios ou formados
apenas por marcação de palco. Os turnos de presidência concentram 42,7\% das palavras dos falantes
presentes em duas ou mais audiências, e boa parte dessa fala consiste em passar a palavra, controlar o tempo e agradecer.
Descartá-los por completo, porém, eliminaria opiniões reais, pois 336 das 2.105 evidências de UDV
estão em turnos de presidência. O turno de presidência mediano tem 28 palavras, e o limite de 50
palavras mantém 2.264 dos 6.266 turnos de presidência, com 89\% de suas palavras, e 332 das 336
evidências. Os turnos mantidos ficam marcados como de presidência.

Após os três filtros restam 13.190 dos 17.264 turnos, cujo texto é copiado da transcrição sem
limpeza; nos 13.190, foi conferido que o texto é idêntico ao intervalo indicado pelos offsets. Os
301 atores que falam em duas ou mais audiências constituem o material dos perfis
(Quadro~\ref{lst:ator}).

\begin{lstlisting}[float=tbp,caption={Exemplo de registro de falas por ator, restrito às audiências
de treino. Os textos foram truncados em {[\ldots]}, e o campo com a concatenação dos turnos de cada
audiência foi omitido.},label=lst:ator]
{
  "actor": "DANIEL BARBOSA",
  "has_party_header": true,
  "party_uf": ["Bloco/PP - AL"],
  "hearings": [
    {"hearing_id": 31, "turns": [
      {"turn_index": 4, "role": "speaker", "start_char": 23770, "end_char": 25429,
       "text": "Obrigado, Deputada Socorro Neri.\n\nAs palavras de V.Exa. foram muito bem colocadas. [...]"}]},
    {"hearing_id": 168, "turns": [
      {"turn_index": 17, "role": "speaker", "start_char": 73694, "end_char": 74807,
       "text": "Bom, Presidente Alice, eu vou fazer uma fala rápida. Eu fui a favor da retirada do Ministério da Ciência, Tecnologia e Inovação do teto de gastos [...]"}]}
  ]
}
\end{lstlisting}

\subsection{Perfis por ator}
\label{sec:metodo-perfis}

A fala de um ator se distribui por audiências de temas diferentes e se mistura a passagens de
condução da sessão e de cortesia. O perfil condensa essa fala em um texto, escrito por um LLM a
partir das falas e destinado a servir de instrução para outro modelo, que deve responder como a
pessoa responderia, tanto sobre temas que ela já tratou quanto sobre temas novos. Por isso, além do
que a pessoa defende, critica, propõe e cobra, o perfil registra os critérios que ela invoca para
justificar essas posições, o que ela declara sobre si e sobre o próprio lado e a forma como argumenta.
Cada perfil é gerado somente com as audiências de treino (Seção~\ref{sec:dados-particao}), para que
nenhuma fala ou opinião das audiências de teste chegue ao prompt.

Cada ator recebe uma única geração, cujo prompt do usuário contém todas as suas falas de treino,
organizadas em blocos por audiência e em ordem cronológica. Cada bloco é aberto por um cabeçalho com a
data da matéria (Seção~\ref{sec:dados-particao}) e o assunto registrado no resumo estruturado, e os turnos de
presidência são precedidos da marca \texttt{[presidência da sessão]}. A data permite situar no tempo
eventuais mudanças de posição, e o assunto ajuda a resolver referências vagas, como ``este projeto
de lei''. O prompt de sistema impõe quatro grupos de regras; os dois prompts estão no
Apêndice~\ref{ap:prompt}.

\textbf{Fonte única.} O perfil deve usar somente as falas, sem conhecimento prévio sobre a pessoa,
os partidos, as organizações ou os temas em debate; por isso, siglas não são expandidas, nomes e
cargos de pessoas citadas não são completados e nenhum fato recebe uma data que a fala não informa.
O cabeçalho de cada bloco vem da matéria e serve para situar as falas, e nada que apareça só nele
pode ser atribuído à pessoa. Dentro das falas, só entra como posição da pessoa o que ela própria
defende: dados citados aparecem como afirmação dela (``afirma que''), e propostas ou opiniões de
terceiros que ela apenas relata, resume ou questiona não entram como posições suas. Perguntas,
hipóteses e condicionais não se tornam afirmações categóricas, e números mantêm o recorte a que se referem. Quando as falas são escassas ou
protocolares, o perfil traz só os blocos que elas sustentam, com um ou dois itens cada.

\textbf{Condução e cortesia.} Em qualquer turno, as passagens de condução e de protocolo, como passar
a palavra, controlar o tempo, citar requerimentos, cumprimentar, agradecer e elogiar, devem ser
desconsideradas, e dos turnos de presidência só se aproveitam os trechos em que a pessoa se
posiciona, argumenta ou propõe.

\textbf{Blocos do perfil.} O perfil é organizado em até quatro blocos de itens, sempre na mesma
ordem, e o bloco sem apoio nas falas é omitido. O bloco de posições registra o que a pessoa defende,
critica, propõe e cobra, e de quem, com o vocabulário do domínio presente nas falas e, quando a
pessoa o declara, com o motivo da posição nas palavras dela. O bloco de critérios e valores reúne os
princípios que a pessoa enuncia de forma geral e os motivos que ela usa para justificar mais de uma
posição, cada um ligado às posições que justifica, e proíbe deduzir um valor a partir de uma posição:
``defende o SUS'' não autoriza ``valoriza o papel do Estado''.

O bloco de alinhamentos declarados contém somente o que a pessoa diz sobre si, como cargo,
organização representada e trajetória, e sobre o próprio lado, como pertencer à base do governo ou à
oposição, sem inferência a partir do partido, das posições ou dos elogios que ela faz. O bloco de
forma de argumentar descreve comportamento observável, como o tipo de apoio usado e o modo de tratar
convidados e colegas, sem adjetivos de personalidade. Um traço só entra nesse bloco se aparecer em
mais de um turno e não for contrariado por outra fala, e a mudança de tom entre audiências é
registrada com a audiência de cada tom, sem ser fundida em um traço único.

\textbf{Escrita.} O perfil é escrito em terceira pessoa, cada bloco abre com um título fixo e cada
item começa pelo verbo (``Defende que'', ``Critica'', ``Declara-se''). Dentro de cada bloco, os
itens são ordenados por importância, com prioridade para o que se repete em mais audiências e ocupa
mais fala, e só levam data quando registram uma mudança de posição entre audiências. Frases sem
conteúdo verificável são cortadas, e o perfil nunca descreve o que a pessoa não disse. O comprimento
é proporcional ao material, com teto de 800 palavras para o conjunto dos blocos; o teto é uma
instrução do prompt e não é imposto durante a geração.

A divisão em blocos separa o material que orienta respostas sobre temas já tratados, as posições, do
material que orienta respostas sobre temas novos, os critérios, os alinhamentos declarados e a forma
de argumentar. As proibições de inferência nesses blocos impedem que o perfil atribua à pessoa
características do seu partido ou das suas posições que ela própria não enunciou; com um alinhamento
deduzido do partido, o modelo que recebe o perfil responderia com base nesse partido, sem apoio nas
falas da pessoa. Pelo mesmo motivo, o perfil não descreve ausências, porque uma frase como ``não
menciona temas de saúde'' poderia ser tomada como posição.

As regras sobre o cabeçalho, sobre posições de terceiros e sobre siglas e nomes foram acrescentadas
após a leitura de perfis gerados durante o desenvolvimento do prompt. Nesses perfis, o assunto da
audiência aparecia como posição do ator, a proposta de um convidado resumida pelo presidente
aparecia como proposta do próprio presidente, e siglas e sobrenomes eram completados com
conhecimento externo às falas. As restrições do bloco de forma de argumentar têm a mesma origem: uma
frase dirigida a uma única autoridade aparecia como traço geral da pessoa, e um traço atribuído a
ela era contrariado por falas de outra audiência.

Os dois prompts compõem uma conversa de duas mensagens, formatada pelo \textit{chat template} do
próprio LLM, sem formatação adicional. A geração usa amostragem com temperatura 0,2, \textit{top-p} de 0,9 e limite de
3.000 tokens de saída, e a semente é fixada antes de cada geração, para que o perfil de um ator não
dependa da ordem em que os atores são processados. Uma resposta que atinge o limite de tokens sem
terminar é considerada falha e descartada. Cada perfil registra o modelo, a versão dos prompts, as
audiências usadas e o número de tokens de entrada e de saída. O modelo gerador é <modelo>.

% Pendente: trecho de um perfil real (Quadro) depois da rodada de geração com o modelo final.

\subsection{Simulação de atores}
\label{sec:metodo-simulacao}

Como o perfil foi escrito para orientar outro modelo a responder como a pessoa responderia, a
simulação o fornece a um LLM, o modelo simulador, e verifica se as respostas reproduzem posições que
a pessoa assumiu em audiências que não entraram no perfil. Três componentes são acrescentados um de
cada vez, para que o efeito de cada um possa ser medido separadamente, o que define quatro condições
encadeadas: a condição 0 recebe só o nome e o cargo da pessoa, a condição 1 acrescenta o perfil, a
condição 2 acrescenta trechos das falas de treino da pessoa recuperados pelo assunto da audiência, e
a condição 3 combina a condição 2 com a condição 0 por \textit{classifier-free guidance} (CFG). Na
execução completa, o modelo simulador é o mesmo LLM que gera os perfis.

\textbf{Prompt comum.} Todas as chamadas partem do mesmo prompt de sistema, e as condições diferem
apenas no material incluído, porque uma instrução diferente em cada condição faria a diferença entre
condições vizinhas, e o prompt negativo da CFG, dependerem também da troca de instrução
(Apêndice~\ref{ap:prompt-simulacao}). O prompt apresenta a pessoa pelo nome e pelo cargo registrado
na matéria sem o partido e a UF, e proíbe o uso de conhecimento prévio sobre a pessoa, os partidos e
o grupo que ela representa, bem como da opinião do próprio modelo sobre o tema. Quando o material não
trata do assunto, o modelo deve se orientar pelos critérios e pelos alinhamentos da pessoa, sem
atribuir a ela posições ausentes do material. O partido é omitido em todas as condições para que a
falta de material não seja suprida pela posição associada ao partido. Cada item do perfil recebe um
identificador formado pela inicial do bloco e pela ordem do item (P1, C1, A1, F1), e na condição 0 o
perfil é substituído por ``Nenhum.'', sem outra alteração nas instruções.

\textbf{Trechos recuperados.} O perfil condensa todas as falas de treino em um texto com teto de 800
palavras e perde detalhes que podem decidir uma pergunta específica. Para recuperá-los, o assunto da
audiência é usado como consulta, codificado pelo mesmo Serafim 335m da UDV, e a busca se restringe
às sentenças de quatro palavras ou mais dos turnos da própria pessoa nas audiências de treino. Entram
as $k$ sentenças de maior cosseno, no máximo uma por turno, e cada trecho mostra a sentença com a
anterior e a seguinte do mesmo turno. Os trechos são ordenados por data, para que uma mudança de
posição apareça na ordem em que ocorreu, recebem identificadores de T1 a T$k$ e trazem a data e o
assunto da audiência de origem, além da marca de presidência quando vêm de um turno de presidência;
o cosseno não é mostrado ao modelo. Os trechos ficam na mensagem do usuário, junto ao pedido, e não no
prompt de sistema, porque variam a cada pergunta, enquanto o perfil é fixo por pessoa.

\textbf{\textit{Classifier-free guidance}.} Mesmo com o perfil no prompt, o modelo pode responder com a
própria opinião sobre o tema ou com o que associa ao nome da pessoa. A CFG, proposta para modelos de
difusão~\citep{ho2022classifier} e adaptada a modelos de linguagem~\citep{sanchez2024stay}, amplifica
a diferença entre um prompt completo $x$ e um prompt negativo $\bar{x}$ a cada passo da geração:
\begin{equation}
  \log \tilde{p}(y_t \mid y_{<t}) = \log p(y_t \mid \bar{x}, y_{<t})
  + \gamma \left[ \log p(y_t \mid x, y_{<t}) - \log p(y_t \mid \bar{x}, y_{<t}) \right],
  \label{eq:cfg}
\end{equation}
em que $\gamma = 1$ reproduz a geração com $x$ e $\gamma > 1$ desloca a distribuição na direção
indicada pelo material presente em $x$ e ausente em $\bar{x}$. O prompt negativo é o da condição 0,
com o perfil ``Nenhum.'', sem trechos e sem exemplos, de modo que a CFG amplifica tudo o que o
material da pessoa acrescenta ao nome e ao cargo. Na múltipla escolha, a condição 3 não exige
passagem adicional pelo modelo, pois aplica a Equação~\ref{eq:cfg} aos log-probs das letras já
calculados nas condições 0 e 2; na geração aberta, a combinação é feita a cada token da fala, o que
exige duas passagens pelo modelo por token.

\textbf{Pedidos.} Na múltipla escolha, usada na avaliação (Seção~\ref{sec:setup}), o modelo indica
qual de quatro opiniões a pessoa emitiu na audiência, e a resposta é lida nos logits das letras A a
D na primeira posição da resposta, sem geração de texto. Os outros três pedidos usam decodificação
gulosa. No pedido da base da estimativa, o modelo classifica o material sobre o assunto como
DIRETA, quando um item de posição ou um trecho trata do mesmo assunto, INDIRETA, quando só há
material sobre assunto vizinho, ESPECULATIVA, quando só critérios ou alinhamentos se aplicam, ou SEM
BASE. Na geração aberta, o modelo escreve em primeira pessoa uma fala de até 200 palavras sobre o
assunto, sem cumprimentos, e, numa chamada separada, a justificação dessa fala. Com SEM BASE, a fala
não é gerada e a saída é ``o material não permite estimar''. A classificação fica numa chamada
própria, e não na chamada da fala, porque o prompt negativo, sem material, favorece a recusa, e com
$\gamma > 1$ a CFG afastaria a fala da recusa justamente quando ela é devida.

Na geração aberta, o prompt de sistema traz também dois exemplos fixos de fala da pessoa, de 60 a
150 palavras em sentenças inteiras, sorteados com semente fixa entre os turnos que não são de
presidência de duas audiências de treino diferentes e iniciados na segunda sentença do turno, porque
os cumprimentos se concentram na primeira. Os exemplos transmitem o estilo da pessoa por fala real, e
não apenas pela descrição do bloco de forma de argumentar, e ficam fora da múltipla escolha, em que a
resposta é uma letra e em que, por conterem posições, eles fariam a diferença entre as condições 0 e
1 medir mais do que o perfil.

\textbf{Justificação conferida.} Na justificação, o modelo recebe a fala simulada e escreve, para
cada frase, um objeto JSON com os identificadores do material que a apoia (itens do perfil, exemplos
e trechos) e, para cada exemplo ou trecho citado, a parte copiada literalmente. A conferência é feita
por código, frase a frase: cada identificador precisa existir no material da chamada, e cada cópia
precisa ter ao menos seis palavras e ser encontrada no exemplo ou trecho correspondente pela regra de
prefixo da citação literal (Seção~\ref{sec:metodo-udv}), que compara as primeiras seis a dez
palavras. Contam como sem fundamento a frase sem apoio, com identificador inexistente ou com cópia
não encontrada, e a frase da fala que não aparece em nenhum objeto. A conferência mostra que o item
citado existe e que o início da cópia é literal, mas não que o item sustenta a frase nem que o modelo
o usou ao escrever a fala, pois a justificação é produzida depois dela; a sustentação de sentido
exige leitura humana.
